\documentclass[11pt]{article}
\usepackage[a4paper, margin=2.5cm]{geometry}
\usepackage[slovene]{babel}
\usepackage[utf8]{inputenc}
\usepackage[T1]{fontenc}
\usepackage{lmodern}
\usepackage{amsmath}
\usepackage{amsthm}
\usepackage{amssymb}
\usepackage{graphicx}

{\theoremstyle{definition}
\newtheorem{definicija}{Definicija}}
{\theoremstyle{plain}
\newtheorem{izrek}{Izrek}}

\title{Brownovo gibanje}
\author{Matej Rojec}
\date{}

\begin{document}
\maketitle

Brownovo gibanje (več v [1]) je intuitivno slučajen proces, \ref{knjiga.bib} % Sklic na knjigo
ki predstavlja naključno gibanje delcev v mediju.
    
    
   \includegraphics{PerrinPlot2.pdf}
     \caption{Reprodukcija slike iz Jean Baptiste Perrin, \emph{Mouvement brownien et réalité moléculaire}, Ann. de Chimie et de Physique (VIII) 18, 5-114, 1909}
     \label{fig:JeanBaptistePerrin}
    
    
    \begin{definicija}
    Standardno Brownovo gibanje \(\{B_t\}_{t \geq 0}\) je slučajen proces z naslednjimi lastnostmi: 
     \begin{itemize}
         \item\(B_0 = 0\).
         \item Prirastki \(B_{t_n} - B_{t_{n-1}}, B_{t_{n-1}} - B_{t_{n-2}}, \ldots, B_{t_2} - B_{t_1}, B_{t_1} - B_{t_0}\) so neodvisne slučajne spremenljivke, za vsak \(t_0 \leq t_1 \leq \cdots \leq t_{n-1} \leq t_n\).
         \item Za vsak \(t \geq 0\) in \(h > 0\) velja \(B_{t+h} - B_t \sim \mathcal{N}(0, h)\).
         \item Funkcija \(t \mapsto B_t\) je zvezna skoraj gotovo.
        \end{itemize}
    \end{definicija}
    
    Preden zapišemo izrek, definirajmo še pojem časa ustavljanja.
    
    \begin{definicija}
    Slučajna spremenljivka \(\tau\) na verjetnostnem prostoru (\Omega, \mathcal{F}, P) z vrednostmi v \mathbb{R}^+
    je čas ustavljanja glede na filtracijo če velja: (\mathcal{F}_t)(t\in T), \forall t\in T:{\(\tau\)\leq t} \in (\mathcal{F}_t).
    \end{definicija}
    
    Zdaj lahko zapišemo izrek \ref{thm:stopped_brownian}. 
    
    \begin{izrek}[Brownovo gibanje]
     \label{thm:stopped_brownian}
    Naj bo \(\{B_t\}_{t \geq 0}\) (standardno) Brownovo gibanje, \(\tau\) čas ustavljanja glede na
    (\mathcal{F}_t)_t\geq 0 in naj velja P[\(\tau\)\)\le \infty].
    Potem je tudi proces:
    \[
    \hat{B} := \{B_{\tau+t} - B_\tau \mid t \geq 0\}
    \]
    (standardno) Brownovo gibanje in neodvisen od (\mathcal{F})_\(\tau\).
    \end{izrek}
    
\bibliographystyle{plain}
\bibliography{knjiga.bib}

\end{document}